\SourceLeaf{288}%
\begin{center}
\begin{tabular}{lr}
Brocken&\(+0''00141\)\\
Hohehagen&\(-0{,}00094\)\\
Inselsberg&\(-0{,}00511\)
\end{tabular}
\end{center}
\Needspace{14\baselineskip}
\hspace*{\parindent}I take a second example from the trigonometric survey of Switzerland\footnote{Ergebnisse der trigonometrischen Vermessungen in der Schweiz, herausgegeben von J. \textsc{Eschmann}. Zurich 1840. S. 79. 99. 189. 190. 196.}, where the largest principal triangle, between Chasseral, Suchet, and Berra, falls right at the limit of the range of our auxiliary table. For these points we have
\begin{center}
\begin{tabular}{l|c|c|r}
&\multicolumn{2}{c|}{Latitude.}&\\
&on the ellipsoid&on the sphere&\(k\)\\\hline
Chasseral&\(47^0\quad8'\quad1''\)&\(47^0\quad6'\quad33''\)&\(6''137\)\\
Suchet&\(46\quad46\quad23\)&\(46\quad44\quad57\)&\(6{,}948\)\\
Berra&\(46\quad40\quad36\)&\(46\quad39\quad11\)&\(7{,}173\)
\end{tabular}
\end{center}
The logarithms of the sides of the triangle in meters are
\begin{center}
\begin{tabular}{lr}
Suchet-Berra&\(4{,}7474503\)\\
Berra-Chasseral&\(4{,}7133766\)\\
Chasseral-Suchet&\(4{,}7808768\)
\end{tabular}
\end{center}
The azimuths
\begin{center}
\begin{tabular}{lrrr}
&\multicolumn{3}{c}{Station Chasseral}\\
Suchet&\(48^0\)&\(36'\)&\(41''\)\\
Berra&349&21&54\\
&\multicolumn{3}{c}{Station Suchet}\\
Chasseral&228&10&40\\
Berra&280&47&19\\
&\multicolumn{3}{c}{Station Berra}\\
Suchet&101&18&40\\
Chasseral&169&27&22
\end{tabular}
\end{center}
\hspace*{\parindent}This gives the reductions of the spheroidal azimuths to the spherical azimuths
